% ====================================================================
\section{Results}
\label{sec:resultsAndDiscussion}
% ====================================================================

\subsection{Experimental Setup}
\label{sec:setup}
The extensions are evaluated with a discrete-round Python simulator of
the IoDT factor model.  It follows the topology of the reference
scenario of~\cite{soula2024iodt}: eight sensors, two of which
(COMP1 and COMP2, labelled C1 and C2 in Fig.~\ref{fig:architecture}) are compromised, attached to two fog nodes, with
twins and the trust engine in the cloud.  The simulator models the
trust logic only; it is not a re-implementation of the iFogSim
evaluation of~\cite{soula2024iodt}, whose resource and energy model it
does not include~\cite{gupta2017ifogsim}.  Delays are counted in
evaluation rounds.  Where a time is given, it assumes a round length
$T_{\mathrm{round}}=117$~ms, obtained by dividing the $1875$-ms
execution time reported for the IoDT scenario in~\cite{soula2024iodt}
into 16 rounds; the
original model itself has no rounds, so this conversion is an
assumption.  A run lasts $T=16$ rounds unless stated otherwise.

In every round each factor of every node is corrupted independently
with probability $q$ (assumption A1).  Because the congestion level of
a real deployment is unknown, results are reported over a range of
$q$.  Two levels are tied to the original evaluation, in which one of
six honest nodes (N3) was misidentified in $2\%$ of
cases~\cite{soula2024iodt}, a figure whose unit is not stated.
\emph{Reading~A} takes it as a per-node, per-round rate of $2\%/6$,
which the binary AND approximately reaches at $q=0.001$ (per-round
FPR $0.40\%$ against $0.33\%$).  \emph{Reading~B} takes it as a
per-run false-isolation rate over the 16-round scenario, again
divided among the six honest nodes ($2\%/6$ per node), which the AND
reaches at $q=5.2\times10^{-5}$.  Neither calibration
reproduces the detection side: the simulated AND detects both attackers
in every run, against $92\%$ and $95\%$ in~\cite{soula2024iodt},
because a lie here always fails $W$ and $T$, while the causes of the
missed detections in the iFogSim runs are not reported.  The stress
level is $q=0.04$.  A healthy PDR is drawn uniformly from $[92,100]\%$
and a degraded one from $[60,89]\%$, against the $90\%$ threshold.

Extension~I is compared with the binary AND of~\cite{soula2024iodt}
and with two delay-matched rules: the AND required to fail in two
consecutive rounds (\emph{debounced AND}), which has the same one-round
delay as Extension~I, and a rule that flags when two of the four
factors fail in the same round or the persistence check fires
(\emph{2-of-4 rule}).  A beta reputation score~\cite{josang2002beta}
on the AND observation, tuned on seed~1 by an analogous rule
(naive liar flagged within two rounds without congestion, lowest FPR
at $q=0.04$, near-ties resolved towards the simpler setting),
reduced to the debounced AND (its FPR was within $0.002$ percentage
points in every cell of Table~\ref{tab:headline}) and is not shown.
The \emph{full system} is Extension~I in each of five replicas with a
$3$-of-$5$ quorum.  Because its replicas receive identical evidence
(B4), its verdicts coincided with those of Extension~I in every cell
of Table~\ref{tab:headline} (its row is therefore omitted), also with two corrupted replicas that
shield the attackers and frame honest nodes;
Section~\ref{sec:res-ext2} evaluates corrupted replicas in detail.

The parameters of Extension~I were selected on a tuning seed (seed~$1$)
by two rules: the lowest per-round FPR at $q=0.04$ among settings that
flag a naive liar within two rounds on average, which selects
$(\alpha,\theta)=(0.5,0.70)$; and the lowest FPR among persistence
windows that flag a persistent minimal-footprint liar within four
rounds and detect one that lies in half of the rounds in at least
$90\%$ of runs, which selects $4$ of $8$.  Every reported cell uses
seed~$0$ with $2{,}000$ runs unless stated otherwise (the long-horizon
runs use seeds $100$--$105$).  Per-round rates carry $95\%$
percentile-bootstrap intervals over runs, which respect the correlation
between rounds of a run, and run-level proportions carry $95\%$ Wilson
intervals.  For a zero count we give the one-sided $95\%$ bound $3/N$
over the $N$ honest node-runs ($0.025\%$ for $N=12{,}000$), which also
bounds the per-round FPR.  Table~\ref{tab:params} lists the defaults.
The experiments ran in Python~3.13 with NumPy and the
\texttt{cryptography} package (OpenSSL) on an Intel Xeon at $2.8$~GHz,
single-threaded.  The tuning grids are provided with the
supplementary material.

\begin{table}[!t]
\caption{Default parameters of the evaluation.}
\label{tab:params}
\centering
\renewcommand{\arraystretch}{1.1}
\footnotesize
\setlength{\tabcolsep}{3pt}
\begin{tabular}{lll}
\toprule
\textbf{Symbol} & \textbf{Meaning} & \textbf{Default} \\
\midrule
$w_E,w_P,w_T,w_W$ & weights (\ref{eq:weighted_score}) & $.3,.3,.2,.2$ \\
$\alpha,\ \theta$ & EMA factor, threshold (\ref{eq:ema}) & $0.5,\ 0.70$ \\
$m,\ w$ & persistence check & $4$ of $8$ rounds \\
$m_2,\ w_2$ & optional second window & off ($8$ of $32$) \\
$q$ & per-factor congestion rate & $5.2\times10^{-5}$--$0.04$ \\
$T$ & horizon of a run & $16$ rounds \\
$R,\ \theta_r$ & re-admission after & $8$ rounds with $\tau_t\ge0.9$ \\
$n,\ k$ & replicas, quorum & $5,\ 3$ \\
$K$ & header (commit) period & $10$ rounds \\
Runs & Monte-Carlo runs per cell & $2{,}000$ \\
\bottomrule
\end{tabular}
\end{table}

% ====================================================================
\subsection{Congestion Tolerance}
\label{sec:res-ext1}
Table~\ref{tab:headline} compares the aggregators at five congestion
levels.  Every aggregator detected both naive attackers in every run at
every level, so the table reports what differs: false positives and
delay.

\begin{table*}[!t]
\caption{Congestion tolerance of Extension~I and the baselines (2{,}000 runs, 16 rounds, two significant figures; 95\% intervals in brackets; zero counts with their one-sided 95\% bound).}
\label{tab:headline}
\centering
\renewcommand{\arraystretch}{1.05}
\setlength{\tabcolsep}{5pt}
\footnotesize
\begin{tabular}{lrrrrr}
\toprule
 & & & \textbf{TTD, naive} & \multicolumn{2}{c}{\textbf{Minimal-footprint liar}} \\
\cmidrule(lr){5-6}
\textbf{Aggregator} & \textbf{Per-round FPR \%} & \textbf{False isolation \%} &
\textbf{liar (rounds)} & \textbf{Det.\ \%} & \textbf{TTD (rounds)} \\
\midrule
%%T:tab_headline%%
\bottomrule
\end{tabular}
\end{table*}

Three regimes appear.  At reading~B the AND falsely isolates $0.39\%$
of honest nodes within 16 rounds, close to the $0.33\%$ that the
reading targets, and every rule that needs more than one failure
isolates none; there is little to improve.  At reading~A the AND
isolates $6.2\%$, while the debounced AND and the 2-of-4 rule isolate
$0.017\%$ and $0.0083\%$ (two and one of $12{,}000$ honest node-runs)
and Extension~I none.  At this level any rule that needs two failures
removes almost all false isolations, and the 2-of-4 rule does so
without delaying detection of a naive liar (a minimal-footprint liar
still takes it about four rounds).  Only as congestion grows does smoothing
make the difference: at $q=0.005$ Extension~I still records no false
flag where the debounced AND isolates $0.56\%$ and the 2-of-4 rule
$0.32\%$; at $q=0.02$ the figures are $0.18\%$, $8.7\%$ and $3.8\%$;
and at $q=0.04$ they are $1.1\%$, $26\%$ and $14\%$, against $93\%$
for the AND.

Fig.~\ref{fig:qsweep} shows why.  The fitted log--log slopes of the
per-round FPR are $0.99$ for the AND ($q$ from $0.0005$ to $0.02$),
$2.0$ for both two-failure rules ($q$ from $0.002$ to $0.04$), and
$3.1$ for Extension~I ($q$ from $0.02$ to $0.08$; below $0.02$ it
records too few flags to fit), in line with the $\Theta(q)$,
$\Theta(q^{2})$ and $\Theta(q^{3})$ orders of
Section~\ref{subsec:ext1-math}.  Equation~(\ref{eq:fpr_and}) predicts
$15.07\%$ for the AND at $q=0.04$, against $15.04\%$ measured.  The
price of smoothing is one round: Extension~I isolates a naive liar after
$1.96$--$2.00$ rounds, like the debounced AND, where the AND and the
2-of-4 rule need one.  Extension~I detects the minimal-footprint liar
in every run at every level, after $3.9$--$4.0$ rounds, as
Proposition~\ref{prop:persist} predicts.  Above $q\approx0.08$ every
aggregator degrades: at $q=0.12$ Extension~I falsely isolates $30\%$
of honest nodes within 16 rounds.

\Figure[!t](topskip=0pt,botskip=0pt,midskip=0pt)[width=\columnwidth]{figs/fig_qsweep.png}
{Per-round FPR (top, logarithmic axes) and false isolation within 16 rounds (bottom) against the per-factor congestion rate, with dotted lines at reading~A and at the stress level.\label{fig:qsweep}}

% ====================================================================
\subsection{Long Operation and Re-admission}
\label{sec:res-long}
A deployed trust engine runs for far longer than 16 rounds.
Table~\ref{tab:long} extends the horizon to
$4{,}096$ rounds ($600$ runs per cell, attackers honest).  False
isolation grows with the horizon for every rule: at $q=0.04$
Extension~I falsely isolates $1.2\%$ of honest nodes within 16 rounds
in these runs ($1.1\%$ in Table~\ref{tab:headline}), $20\%$ within
$256$ and $60\%$ within $1{,}024$, while the debounced AND and the
2-of-4 rule reach $99\%$ and $91\%$ within $256$ rounds.

The mean time to the first false isolation (MTTFI) does not depend on
the horizon.  At reading~A it is $250$ rounds for the AND (about
$29$~s at the assumed period), $6.2\times10^{4}$ rounds for the
debounced AND ($2.0$~h) and $1.5\times10^{5}$ for the 2-of-4 rule
($5.0$~h).  Extension~I recorded no flag in $1.5\times10^{7}$ honest
node-rounds, so its MTTFI is at least $4.9\times10^{6}$ rounds ($6.7$
days) at $95\%$ confidence.  At $q=0.005$ its MTTFI is
$7.4\times10^{5}$ rounds ($24$~h), against $2.6\times10^{3}$ rounds
($5$~min) and $6.6\times10^{3}$ rounds ($13$~min) for the two simple
rules; at $q=0.04$ it is $1.1\times10^{3}$ rounds ($2.2$~min), against
$52$ and $110$ rounds.

\begin{table}[!t]
\caption{Long operation (600 runs of 4{,}096 rounds; no attackers): false isolation within 1{,}024 rounds, MTTFI, and quarantine with re-admission.}
\label{tab:long}
\centering
\renewcommand{\arraystretch}{1.08}
\footnotesize
\setlength{\tabcolsep}{2.4pt}
\begin{tabular}{lrrrr}
\toprule
 & \textbf{False isol.} & \textbf{MTTFI} & \textbf{Quarant.} & \textbf{False quar.} \\
\textbf{Aggregator} & \textbf{1{,}024 rds \%} & \textbf{(rounds)} & \textbf{time \%} & \textbf{/node/h} \\
\midrule
%%T:tab_long%%
\bottomrule
\end{tabular}
\end{table}

With quarantine and re-admission (Section~\ref{subsec:metrics}), a
falsely quarantined honest node returns after $8$ to $18$ rounds on
average ($9$ to $16$ under Extension~I).  At $q=0.04$ an honest node
then spends $1.4\%$ of its time in quarantine under Extension~I ($27$
false quarantines per node and hour), against $7.2\%$, $15\%$ and $73\%$
under the 2-of-4 rule, the debounced AND and the AND.  At $q=0.005$ the
share is $0.0013\%$ under Extension~I and $0.12$--$15\%$ under the
others.  Under sustained congestion as high as $q=0.04$, false
quarantines are an operating cost of every rule, which Extension~I cuts
by a factor of $5$ to $50$; at lower congestion it removes them.

% ====================================================================
\subsection{Stealthy, Minimal and Rule-Aware Attackers}
\label{sec:res-stealth}
Table~\ref{tab:stealth} varies the lying rate $p$ of attackers that lie
at random.  Stealth has a price for every rule except the AND, which
pays for its detection with false isolation.  At reading~A,
Extension~I detects a liar that fails $T$ and $W$ in $97\%$ of runs at
$p=0.5$ but in $57\%$ at $p=0.25$ and $13\%$ at $p=0.10$ within
16 rounds, much like the debounced AND; over 64 rounds the last two
figures rise to $97\%$ and $46\%$.  The 2-of-4 rule flags every lie
that fails two factors and detects such a liar almost as well as the
AND ($82\%$ of runs at $p=0.10$), at the false-isolation cost of
Table~\ref{tab:headline}.  A liar that fails only $W$ is
detected by Extension~I in $92\%$ of runs at $p=0.5$ and $32\%$ at
$p=0.25$ ($86\%$ over 64 rounds), about as often as by the 2-of-4
rule, whose persistence check is the same.  At $q=0.04$ every rule detects more, partly because
congestion adds failures to the attacker's own record.

\begin{table}[!t]
\caption{Detection within 16 rounds (\%) of attackers that lie at random with probability $p$ (2{,}000 runs; the false isolation of each rule is that of Table~\ref{tab:headline}).}
\label{tab:stealth}
\centering
\renewcommand{\arraystretch}{1.08}
\footnotesize
\setlength{\tabcolsep}{1.6pt}
\begin{tabular}{lrrrrrrrr}
\toprule
 & \multicolumn{4}{c}{\textbf{Lie fails $T$ and $W$}} & \multicolumn{4}{c}{\textbf{Lie fails $W$ only}} \\
\cmidrule(lr){2-5}\cmidrule(lr){6-9}
$p$ & AND & Deb. & 2-of-4 & Ext.~I & AND & Deb. & 2-of-4 & Ext.~I \\
\midrule
%%T:tab_stealth%%
\bottomrule
\end{tabular}
\end{table}

A rule-aware attacker does better (Table~\ref{tab:aware}).  The
simulated attacker predicts its own EMA from its past evidence
(assuming a $96\%$ PDR), keeps a margin of $0.02$ above $\theta$, and
respects every persistence window.  Without congestion it lied in $3$ of
$8$ rounds ($37.5\%$) against Extension~I and in $7$ of $32$ ($21.9\%$)
with a second window of $8$ failures in $32$ rounds, without being
flagged --- the budgets of Proposition~\ref{prop:persist}.  Every lie
still fails $W$, so the AND flags the first one, at a cost of $6.2\%$
false isolation at reading~A.  At reading~A congestion rarely adds
failures: over 256 rounds Extension~I detects the attacker in $37\%$ of
runs ($T$ and $W$) and $20\%$ ($W$ only), and with re-admission the
attacker delivers a lie in $36\%$ of rounds ($22\%$ with the second
window).  At $q=0.04$ congestion pushes the attacker over a trigger more
often, and the delivered share falls to $19\%$ ($14\%$); a rule-aware
liar that fails only $W$ delivers $25\%$ ($16\%$).  The simulated
attacker is tuned to Extension~I and faces every defence with the
same strategy, so the AND and debounced-AND figures show what that
attacker suffers, not what an attacker tuned to them would achieve;
against the debounced AND it delivered $16\%$ of rounds ($T$ and $W$)
and $4.5\%$ ($W$ only) at $q=0.04$.  Analytically, the simple rules
leave budgets of their own: the debounced AND never sees lies that
are one round apart (budget $50\%$), and the 2-of-4 rule flags every
lie that fails $T$ and $W$ but leaves a $W$-only liar the same
$37.5\%$ budget.  Threshold rules with published parameters thus leave
a lie budget that a second, longer window shrinks but does not remove.

\begin{table}[!t]
\caption{Rule-aware on-off attacker tuned to Extension~I (all defences face the same strategy): lie budget without congestion, detection within 256 rounds at reading~A (500 runs), lies delivered per round with re-admission (400 runs of 2{,}048 rounds, $T$ and $W$ fail), and the honest nodes' quarantine time at $q=0.04$.}
\label{tab:aware}
\centering
\renewcommand{\arraystretch}{1.08}
\footnotesize
\setlength{\tabcolsep}{2.0pt}
\begin{tabular}{lrrrrrrr}
\toprule
 & \multicolumn{2}{c}{\textbf{Budget \%}} & \multicolumn{2}{c}{\textbf{Det.\ \%}} & \multicolumn{2}{c}{\textbf{Delivered \%}} & \textbf{Hon.} \\
\cmidrule(lr){2-3}\cmidrule(lr){4-5}\cmidrule(lr){6-7}
\textbf{Defence} & $TW$ & $W$ & $TW$ & $W$ & $q{=}.001$ & $q{=}.04$ & \textbf{quar.\,\%} \\
\midrule
%%T:tab_aware%%
\bottomrule
\end{tabular}
\end{table}

% ====================================================================
\subsection{Corrupted, Withholding and Slow Twins}
\label{sec:res-ext2}
Table~\ref{tab:collusion} turns to Extension~II at $q=0.04$, with
identical evidence for all replicas (B4).  A
corrupted single twin hides its node in every run.  With $3$-of-$5$
voting, up to $k-1=2$ corrupted replicas that shield the attackers and
frame honest nodes change nothing: detection stays at $100\%$ of runs
and of steady-state rounds, and the FPR and false isolation equal those
of an honest single twin, as Section~\ref{sec:ext2} argues.  A third
corrupted replica ($c=k$) defeats the vote: the attackers are never
flagged and every honest node is isolated.  Deployments that expect
three replicas at risk need $4$-of-$7$, which tolerates them.

Missing votes are handled by the quorum rule.  With two withholding
replicas plus $5\%$ vote timeouts, or with $20\%$ timeouts, it keeps
false isolation at the level of an honest single twin and detection at
$100\%$ of runs; the cost is rounds with no quorum ($14\%$ and
$5.8\%$), in which verdicts are held and an alarm is raised.  With
$20\%$ timeouts and two withholding replicas, half of all rounds have
no quorum and per-round detection falls to $97\%$, but no missing
vote caused a new isolation.  These timeout rows relax (B3); a node
that is already isolated stays isolated, and cannot be re-admitted,
while no quorum forms.

(B4) matters.  When each replica instead misreads each factor
independently with probability $\eta$ on top of the shared congestion
($q=0.04$), five honest replicas lower false isolation from $12\%$
(one twin) to $2.4\%$ at $\eta=0.05$, but replicas that frame honest
nodes now tip the nodes whose honest replicas disagree: false
isolation rises to $6.0\%$ with one framing replica and $29\%$ with
two, and to $72\%$ with two at $\eta=0.10$ (against $9.3\%$ with
five honest replicas).  Shielding costs little: with two shielding
replicas, steady-state detection falls from $99.9\%$ to $98\%$ at
$\eta=0.05$.  The ``honest twin'' guarantee of Extension~II therefore
holds only when the replicas see the same evidence.

\begin{table}[!t]
\caption{Corrupted, withholding and slow replicas at $q=0.04$, identical evidence (2{,}000 runs; per-round detection in steady state).}
\label{tab:collusion}
\centering
\renewcommand{\arraystretch}{1.08}
\footnotesize
\setlength{\tabcolsep}{2.0pt}
\begin{tabular}{p{0.40\columnwidth}rrrrr}
\toprule
 & \textbf{Det.} & \textbf{Per-rd} & \textbf{FPR} & \textbf{False} & \textbf{No} \\
\textbf{Setting} & \textbf{\%} & \textbf{det.\,\%} & \textbf{\%} & \textbf{isol.\,\%} & \textbf{quor.\,\%} \\
\midrule
%%T:tab_collusion%%
\bottomrule
\end{tabular}
\end{table}

% ====================================================================
\subsection{Audit Log: Cost and Tamper Evidence}
\label{sec:res-ext3}
Table~\ref{tab:chain} reports the audit log for $100$ sensors over
$30$ rounds of simulator output, with one corrupted replica and $5\%$
vote timeouts.  A header with three endorsements occupies $367$ to
$371$~bytes (DER signatures vary slightly in length) whatever the
number of nodes, so header storage falls from $11.0$~kB at $K=1$ to
$1.1$~kB at $K=10$; at the assumed period and $K=10$ this is about
$27$~MB a day for any fleet size.  Off-chain, the events, signed vote
vectors and state records of the 30 rounds occupy $115$--$152$~kB, of
which $87$~kB are events.

\begin{table}[!t]
\caption{Audit-log storage for 100 sensors over 30 rounds.}
\label{tab:chain}
\centering
\renewcommand{\arraystretch}{1.08}
\footnotesize
\setlength{\tabcolsep}{4pt}
\begin{tabular}{rrrrr}
\toprule
$K$ & \textbf{Headers} & \textbf{Header bytes} & \textbf{Off-chain (B)} & \textbf{Max header (B)} \\
\midrule
%%T:tab_chain%%
\bottomrule
\end{tabular}
\end{table}

Table~\ref{tab:tamper} reports $30$ random edits per cell at $K=10$.
The verifier rejected every edit when the adversary held fewer than
$k$ keys, whether it left the hashes untouched, recomputed them, or
re-signed with $k-1$ stolen keys.  Substituting interior-node values
for events and replaying a vote were rejected even with $k$ keys,
because record lengths and vote bindings are checked.  The other
$k$-key rows measure our implemented forger, not a limit of the
scheme: by Theorem~\ref{thm:tamper} a forger that re-signs every
affected record passes whenever no newer checkpoint exists.  Our
forger duplicated events and truncated windows without re-signing the
affected counts and votes, so it was rejected at $K=10$, whereas at
$K=1$ a truncation passed without a checkpoint.  It rewrote a trust
value, a state record or a vote and passed without a checkpoint or
with an older one, and was always rejected against a newer one.  Its
verdict edits were rejected without a checkpoint in $10\%$ of trials,
all of them rounds in which a stolen replica's vote had timed out and
the forger did not sign a replacement.  Finally, in $30$ trials a trust engine
that wrote a verdict its votes did not support never obtained the $k$
endorsements needed to commit the window.  These trials check the
implementation; the guarantee itself rests on the cryptographic
assumptions.

\begin{table}[!t]
\caption{Edits rejected by the verifier (\%, 30 per cell, $K=10$; the ``older'' column uses the 20--21 trials whose edit follows an earlier header) by edit and adversary; ``--'': not run.}
\label{tab:tamper}
\centering
\renewcommand{\arraystretch}{1.08}
\footnotesize
\setlength{\tabcolsep}{2.2pt}
\begin{tabular}{lrrrrrr}
\toprule
 & \multicolumn{2}{c}{\textbf{No keys}} & \textbf{$k{-}1$} & \multicolumn{3}{c}{\textbf{$k$ keys, checkpoint}} \\
\cmidrule(lr){2-3}\cmidrule(lr){5-7}
\textbf{Edit} & \textbf{edit} & \textbf{rehash} & \textbf{keys} & \textbf{none} & \textbf{older} & \textbf{newer} \\
\midrule
%%T:tab_tamper%%
\bottomrule
\end{tabular}
\end{table}

% ====================================================================
\subsection{Robustness}
\label{sec:res-robust}
Fig.~\ref{fig:robust} relaxes assumptions (A1) and (A2) in four ways.

\emph{Bursts.}  Holding the mean congestion at $q=0.04$, a two-state
Markov chain alternates between $q=0.02$ and $q=0.20$ (one ninth of
the time in the high state) with mean bursts of $1$ to $32$ rounds.
Longer bursts raise Extension~I's false isolation from $3.7\%$ to
about $9\%$ when each node has its own chain (panel a), and from
$4.2\%$ to about $9.5\%$ when one chain governs the whole network.  The AND
($77$--$91\%$), the debounced AND ($18$--$29\%$) and the 2-of-4 rule
($16$--$32\%$) stay far higher in either case.

\emph{Common-mode events.}  An event that knocks out all factors of
every node behind a fog node, with probability $q_{\mathrm{corr}}$ per
round, cannot be told from an attack on all those nodes at once, and
Extension~I's false isolation rises to $8.4\%$ at
$q_{\mathrm{corr}}=0.005$ and $28\%$ at $0.02$ (panel b).  Because such
events hit many nodes together, we evaluated \emph{domain
suppression}: when more than half of a fog domain is flagged in the
same round, those flags are discarded for that round (an alarm would
be raised; it is not modelled).  It cuts
false isolation to $1.6\%$ and $3.5\%$ at the same rates and left the
detection of both attackers at $100\%$ of runs.  An adversary able to provoke
common-mode events could use suppression to mask its own lies during
them; that case is not evaluated.

\emph{Cross-factor events.}  When one node-level event fails $P$, $T$
and $W$ together ($4{,}000$ runs per point, no other noise), (A1) fails
across factors and Extension~I errs at $\Theta(q^{2})$ (panel c): its
FPR grows from $0.011\%$ to $0.055\%$ to $0.22\%$ as the event rate
doubles from $0.01$ to $0.04$.  The debounced AND is then better
at every rate ($0.15\%$ against $0.22\%$ at $0.04$), and the 2-of-4 rule (not plotted) is as bad
as the AND, since every event fails three factors.  Such an event also
looks exactly like a lie that fails $T$ and $W$ with a degraded PDR.

\emph{Uneven congestion.}  Concentrating the same mean congestion on
one factor ($0.13$ on it, $0.01$ on the others) leaves the AND
almost unchanged ($15.0\%$ to $15.6\%$ per round) but raises Extension~I's per-round FPR eight- to ninefold, to
$0.81$--$0.88\%$, and its false isolation to $4.8$--$5.7\%$ (panel d),
because one factor now fails repeatedly and the persistence check
fires.  Extension~I stays below the 2-of-4 rule ($11\%$), but a
deployment that knows its noisiest factor should set $m$ per factor.

\Figure[!t](topskip=0pt,botskip=0pt,midskip=0pt)[width=\textwidth]{figs/fig_robust.png}
{Robustness to bursty (a), common-mode (b), cross-factor (c) and uneven (d) congestion.\label{fig:robust}}

\emph{Scale.}  With $100$ sensors behind five fog nodes ($500$ runs, $q=0.04$),
Extension~I isolated $1.1\%$ of honest nodes, as in the eight-node
case, and with common-mode events at $q_{\mathrm{corr}}=0.005$ domain
suppression brought false isolation from $9.0\%$ back to $2.0\%$.

% ====================================================================
\subsection{Trace-Driven Evaluation on Public Datasets}
\label{sec:res-trace}
The congestion of the previous sections is independent across factors,
rounds and nodes.  To test the extensions on real loss patterns, we
replaced it, for every node, with the packet-delivery record of a real
deployment.  The Intel Lab trace~\cite{bodik2004intel} records which
readings of $54$ Mica2Dot motes reached a base station, one per mote
every $31$~s, over $36$ days of a lossy multi-hop network ($51\%$ of
readings arrived; timestamp jitter is not counted as loss, and periods
in which almost no mote was heard are logging gaps and are skipped).
The FlockLab dataset~\cite{jacob2019flocklab} records, about every two
hours in August 2019, which of $100$ back-to-back packets each of the
$27$ TelosB nodes of the FlockLab testbed~\cite{lim2013flocklab}
received from every other.  We took as fog gateways the two nodes that
put the most other nodes in range, and let each of the $24$ remaining
nodes in range report to the gateway it hears best (single hop, mean
packet reception $95\%$).  Two of these nodes ran only during the first
day and a half, so FlockLab windows use $22$ nodes and span the whole
month.

Each factor's evidence of a node in a round is one message, sent in up
to $r$ attempts that use consecutive transmission opportunities of
that node in the trace ($31$-s slots for Intel Lab, packets of one
test for FlockLab); the factor fails when all $r$ attempts are lost.
The retry count $r$ thus sets the loss level, which is then shaped
entirely by the trace.  The two traces differ in time scale: an Intel
Lab round lasts $4r$ slots ($8$ to $33$ minutes), whereas a FlockLab
round spans a few milliseconds of back-to-back packets and windows that
cross a test boundary join tests about two hours apart.  A round in
which no evidence of a node arrives at all is treated as
unavailability, and every rule keeps its previous verdict.  Each run
draws a window of $16$ rounds and $32$ (Intel Lab) or $22$ (FlockLab)
nodes alive throughout, two of which are attackers whose lies are added
to their own real losses ($1{,}000$ runs per cell, drawn from $73$ to
$309$ distinct windows).  No parameter was retuned.  Windows start on a
grid of $16$ rounds, and the intervals below resample whole windows.

\begin{table}[!t]
\caption{False isolation within 16 rounds (\%) on two public packet-delivery traces by retry count $r$, with $q$ the share of lost factor messages in the evaluated windows (1{,}000 trace runs; 2{,}000 runs for the i.i.d.\ column).}
\label{tab:trace}
\centering
\renewcommand{\arraystretch}{1.08}
\footnotesize
\setlength{\tabcolsep}{2.4pt}
\begin{tabular}{lrrrrrrrr}
\toprule
 & & & & & & & \multicolumn{2}{c}{\textbf{Ext.~I}} \\
\cmidrule(lr){8-9}
\textbf{Trace} & $r$ & $q$\,\% & AND & Deb. & 2-of-4 & Ext.~I & typ. & i.i.d. \\
\midrule
%%T:tab_trace%%
\bottomrule
\end{tabular}
\end{table}

Table~\ref{tab:trace} gives the result.  Extension~I has the lowest
false isolation in every cell, and the window-level interval of its
difference to each other rule excludes zero in every cell (the
smallest is $0.16$ to $1.1$ percentage points, against the debounced
AND at FlockLab with $r=4$).  The traces narrow its lead, however.  At
Intel Lab with $r=8$ ($7.9\%$ of messages lost) it falsely isolates
$15\%$ [$12$, $19$] of honest nodes within 16 rounds, against
$22$--$24\%$ for the two-failure rules and $42\%$ for the AND; with
$r=16$ the figures are $6.5\%$ [$4.2$, $9.1$], $10$--$11\%$ and $21\%$.
At FlockLab without retransmission it isolates $8.3\%$ [$7.5$, $9.2$]
against $22$--$30\%$ and $48\%$, and with $r=4$ $2.0\%$ [$1.4$, $2.7$]
against $2.6$--$3.1\%$ and $8.0\%$.  Independent congestion at the same
mean rate (column ``i.i.d.'') would make the AND isolate far more nodes
everywhere and Extension~I fewer everywhere except at the highest loss
level (Intel Lab, $r=4$).  Two properties of real loss explain this.
It is concentrated: at Intel Lab with $r=8$ the median node loses
$5.6\%$ of its messages over the whole trace and the lossiest tenth
$19\%$, so the AND's errors fall on a minority of nodes.  And it is
bursty across factors and rounds: in the evaluated windows two or more
factors of a node fail in the same round $1.4$ to $9.2$ times as often
as independent losses would make them, and a failed factor fails again
in the next round in $42$--$76\%$ of cases, which is the cross-factor
case of Section~\ref{sec:res-robust}, in which smoothing errs at
$\Theta(q^{2})$.  For the nodes whose long-run loss is at most the
median (column ``typ.'', about $10{,}000$ to $15{,}500$ node-runs per
cell), Extension~I's false isolation is $6.6\%$ and $0.91\%$ at Intel
Lab ($r=8$, $16$) and below $0.7\%$ at FlockLab; $65$--$99\%$ of its
false isolations fall on nodes with above-median loss.

Detection is slightly lower than in the synthetic study.  Extension~I
flagged a naive liar in $98.6$--$99.9\%$ of runs after about two rounds,
and a minimal-footprint liar in $98.8$--$99.9\%$ after $3.7$--$4.1$
rounds.  The misses are attacker nodes whose evidence was lost in
almost every round of the window, so that their verdict stayed at
\textsc{trust} under the silence rule; counting silent rounds as
failures instead detects every attacker and raises false isolation by
at most $2.4$ percentage points, with Extension~I still the lowest in
every cell.

Over longer windows false isolation grows for every rule: within $128$
rounds (Intel Lab, $r=8$) and $256$ rounds (FlockLab, $r=1$) Extension~I
isolates $30\%$ and $40\%$ of honest nodes against $41$--$66\%$ and
$61$--$81\%$ for the other rules, and with re-admission honest nodes
spend $14\%$ and $11\%$ of their time in quarantine against
$15$--$37\%$.  The rule-aware attacker, still tuned to Extension~I,
delivered a lie in $28\%$ (Intel Lab) and $27\%$ (FlockLab) of rounds
against it, $17\%$ with the second window, and was detected in
$61$--$74\%$ of runs; the AND and the 2-of-4 rule held it to about
$1.5\%$ of rounds at the cost of the false isolation above.  The
traces thus confirm the ordering of the synthetic study but not its
magnitude: on real, lossy links every rule isolates honest nodes far
more often, and per-node baselines that separate a weak link from
misbehaviour are needed before isolation is automated.

% ====================================================================
\subsection{Latency}
\label{sec:res-latency}
Table~\ref{tab:latency} gives the measured cryptographic cost per round.
Because each replica signs one vote vector per round, the engine
verifies five signatures, $2.0$~ms for $8$ or $100$ nodes, and each
replica signs once ($0.5$~ms); verifying one vote per node would cost the
engine $203$~ms per round at $100$ nodes, more than the assumed
$117$-ms period.  One EMA update for one node costs
$0.1$~$\mu$s; the persistence check was not timed separately.  The remaining term is the time
$d_{\mathrm{rep}}$ for votes to reach the engine, which depends on
replica placement and was not measured: replicas in different
availability zones of one cloud region are a single-digit number of
milliseconds apart~\cite{aws2024az}, whereas replicas at different fog
sites add a fog-to-cloud link ($100$~ms in the configuration
of~\cite{soula2024iodt}) unless the engine runs at the fog.
Extension~III adds nothing to the decision path.  Relative to the AND,
isolating a naive liar therefore takes one more round of evidence plus
about $2.5$~ms plus $d_{\mathrm{rep}}$.  All the added work runs in the
fog and cloud, so the extensions add no computation or energy on the
sensors.

\begin{table}[!t]
\caption{Measured per-round cost on the decision path (ms, $n=5$ replicas, single thread).}
\label{tab:latency}
\centering
\renewcommand{\arraystretch}{1.08}
\footnotesize
\setlength{\tabcolsep}{3pt}
\begin{tabular}{lrr}
\toprule
\textbf{Operation} & \textbf{8 nodes} & \textbf{100 nodes} \\
\midrule
%%T:tab_latency%%
\bottomrule
\end{tabular}
\end{table}

% ====================================================================
\section{Discussion}
\label{sec:discussion}
% ====================================================================

% ====================================================================
\subsection{Comparison With Related Frameworks}
\label{sec:res-comparison}
Table~\ref{tab:comparison} positions the proposed system against the
trust frameworks closest to it.  Each cell reflects what the cited
paper describes, and the latency row reports what each paper measures.
Temporal smoothing is marked $\checkmark$ for an explicit forgetting
factor and $\circ$ for trust that accumulates history without one.
The latency figures are heterogeneous (trust set-up, block operations,
detect-and-isolate time) and come from different platforms, so they
show where each system places its costs rather than ranking the
systems.  Within this set, the proposed design is the only one that
combines a continuous, smoothed trust score, replicated twins with
voting, and a tamper-evident record kept off the decision path.

\begin{table*}[!t]
\caption{Comparison with related trust frameworks ($\checkmark$ yes, $\circ$ partial or optional, -- no).}
\label{tab:comparison}
\centering
\renewcommand{\arraystretch}{1.12}
\setlength{\tabcolsep}{6pt}
\footnotesize
\begin{tabular}{lccccccc}
\toprule
\textbf{Capability} & \textbf{BPT}~\cite{dewanta2021bpt} & \textbf{Debe}~\cite{debe2019iot}
 & \textbf{Hussain}~\cite{hussain2020context} & \textbf{Onwubiko}~\cite{onwubiko2023enabling}
 & \textbf{TMIoDT}~\cite{akram2024tmiodt} & \textbf{IoDT}~\cite{soula2024iodt} & \textbf{Proposed} \\
\midrule
Digital twin              & -- & -- & -- & \checkmark & \checkmark & \checkmark & \checkmark \\
Continuous trust score    & \checkmark & \checkmark & \checkmark & -- & $\circ$ & -- & \checkmark \\
Temporal smoothing        & $\circ$ & $\circ$ & \checkmark & -- & $\circ$ & -- & \checkmark \\
Replicated twins, voting  & -- & -- & -- & -- & -- & -- & \checkmark \\
Blockchain                & \checkmark & \checkmark & -- & \checkmark & \checkmark & -- & $\circ$ \\
Chain off decision path   & $\circ$ & -- & n/a & -- & -- & n/a & \checkmark \\
Tamper-evident record     & \checkmark & \checkmark & -- & \checkmark & \checkmark & -- & \checkmark \\
\midrule
Reported latency          & 2.47\,s & n.r. & n.r. & n.r. & 15--16\,ms & 250--300\,ms & +1 round \\
(what it measures)        & trust set-up & & & (gas cost only) & block op. & detect + isolate & +2.5\,ms\,+\,$d_{\mathrm{rep}}$ \\
\bottomrule
\multicolumn{8}{l}{n.r.\ = not reported.} \\
\end{tabular}
\end{table*}

TMIoDT~\cite{akram2024tmiodt} is the closest precedent.  It assigns
each drone a blockchain-recorded credit score and validates
transactions by proof-of-authority consensus, so the chain is involved
in its trust decisions, and it detects intrusions with deep models trained
on labelled traffic (about $99.1\%$ accuracy on ToN-IoT); its twins are
not replicated.  The proposed score needs no training data and keeps
the log off the decision path, but sees only the four behavioural
factors, so the two approaches are complementary.

\subsection{Limitations}
Behaviour-based trust has hard limits, which the experiments make
explicit.  An attacker that knows the rule keeps a lie budget
($37.5\%$ of rounds against Extension~I without congestion,
$21.9\%$ with a second window), low-rate random liars are detected
slowly, and one that never deviates on any factor cannot be seen at
this layer at all.  Bad-mouthing through the notification factor is
not covered.  Common-mode events need domain suppression, which an
adversary able to provoke them could abuse, and events that fail
several factors at once lower the order of the error to
$\Theta(q^{2})$.  On the architecture side, the fog relay is trusted
(it enforces verdicts and logs the verdicts it enforced), although a
replica hosted at a fog site may be corrupted, and the trust engine is
trusted for timeliness; its integrity is checked after the fact,
within $K$ rounds.  Extension~II assumes identical evidence and
bounded vote delay; with independent replica errors, two framing
replicas raise false isolation several-fold.  An adversary that controls
$k$ replicas defeats any quorum by construction.  Finally, the
evaluation is a simulation of the trust logic; honest losses come from
synthetic congestion or from two public delivery traces mapped onto the
factors by our retransmission model, attackers are simulated, and on
those traces lossy nodes are falsely isolated often under every rule;
under the silence rule a node that goes completely silent keeps its
last verdict, so silence itself must be monitored separately.
The round length, the header append
latency and the replica delay $d_{\mathrm{rep}}$ are assumptions or
were not measured, and no ledger was deployed.
